\scanpage{9}
\mspage{6}
\[
u^4c'_4=c_4\qquad u^6c'_6=c_6\qquad u^{12}\Delta'=\Delta\qquad .\tag*{(1.8)}
\]

Les formules (1.8) expriment que la section $c_4\pi^{\otimes4}$ (resp. $c_6\pi^{\otimes6}$, resp. $\Delta\pi^{\otimes12}$) du fibré $\omega^{\unsure{\tikz[baseline=.06ex,x=1ex,y=1ex]{\draw[line width=.055ex] (.4,.4) circle (.33); \draw[line width=.055ex] (.17,.63)--(.63,.17);}|\otimes}4}$ (resp. $\omega^{\otimes6}$, $\omega^{\bigcirc12}$) ne dépend pas du choix arbitraire de $\pi$, $x$ et $y$.

2.\qquad \uline{Cas où 2 et 3 sont inversibles}.

Lorsque 2 et 3 sont inversibles, on peut d'une et d'une seule façon choisir $x$ et $y$ ($\pi$ étant donné) de sorte que dans l'équation (1.3.) on ait $a_1=a_2=a_3=0$. On a alors
\[
y^2=x^3+a_4x+a_6\qquad,\ \text{avec}\quad\pi=dx/2y\qquad,\tag*{(2.1)}
\]
\[
c_4=-2^4\cdot3a_4,\quad c_6=-2^5\cdot3^3\,a_6,\quad\Delta=2^4(-4a_4^3-27a_6^2)\quad .\tag*{(2.2)}
\]
Posant $Y=2y$, $X=x$, on retombe sur la forme de Weierstrass
\[
Y^2=4X^3-g_2X-g_3\quad,\ \text{avec}\quad\pi=d\,X/Y\quad,\ \text{et}\tag*{(2.3)}
\]
\[
g_2=-4a_4=\frac{1}{12}c_4\qquad g_3=-4a_6=\frac{1}{216}c_6\tag*{(2.4)}
\]
\[
\Delta=g_2^3-27g_3^2\qquad .
\]

\uline{Proposition} 2.5.\qquad \uline{Au dessus de} $\mathrm{Spec}(\mathbb Z[2^{-1},3^{-1}]$, \uline{il existe un schéma de modules pour les courbes de genre un munies d'une forme différentielle invariante inversible}. \uline{C\unsure{;|e} schéma est}
\[
\mathrm{Spec}(\mathbb Z[2^{-1},3^{-1}][g_2,g_3]\qquad ,
\]
% archive sheet number: 8
